% ---- Chris: Bayesian surface habitability mapping ----
\subsection{Surface habitability mapping (Chris)}

My study asks a spatial question that the other four do not: \emph{where} is Titan's surface most likely
to be habitable, and how has that pattern changed through time? Working at the outermost layer of
Figure~\ref{fig:crosssection}, I cast this as Bayesian inference, computing a posterior habitability
probability $P(H)$ for every pixel of a global equirectangular grid ($\sim$$6.5\times10^{6}$ pixels at
4490\,m per pixel, for Titan's radius of 2575\,km). Eight geomorphological and compositional features
-- including liquid-hydrocarbon coverage, organic (VIMS) spectral signatures, geomorphological
diversity, and a methane-cycle proxy -- update a beta-conjugate prior, so that each map carries a
calibrated uncertainty rather than a single score \citep{Meadows2026}; surface temperatures follow
\citet{Jennings2019}.

The framework is applied at five epochs, from the Late Heavy Bombardment ($-3.5$\,Gya) to a red-giant
future ($+5.9$\,Gya), turning a static map into a trajectory. At present the north-polar lake shores
dominate the ranking: the liquid-hydrocarbon feature raises both the maximum and the spread of $P(H)$,
and the north-polar-to-equatorial contrast is the largest of any epoch (Figure~\ref{fig:habmap}). Selk
crater -- the Dragonfly landing site -- in fact scores only modestly ($P(H)\approx0.21$, against
$\sim0.55$ for the leading polar lakes), because the map measures present-day \emph{solvent}
habitability whereas Dragonfly targets \emph{past} liquid-water chemistry; the framework instead
endorses the mission through Selk's geomorphological diversity, the highest of any equatorial site and a
proxy for the organic--water-ice contact terrain it will sample \citep{Lorenz2021}. The global median
$P(H)$ climbs from $0.27$ at the LHB to $0.32$ at present, holds through the near-future plateau
($r>0.999$), then -- after a $\sim$1\,Gyr solvent-free minimum around $+5$\,Gya, once the hydrocarbon
seas evaporate -- rises to its maximum of $0.53$ under a red-giant water--ammonia ocean. Habitability is
therefore episodic in time; spatially, the lowest scores fall on the water-ice-rich cores of Xanadu and
the mountain ranges.

A deliberate strength is that the map is \emph{ocean-agnostic}: it scores the surface on its own
evidence, and so remains valid whether or not Titan retains the contested subsurface ocean
\citep{Petricca2025}, making it a robust anchor for the group's argument. It also indexes the other
studies: my top sites are polar lake shores, where present-day solvent and organics coincide, whereas
the impact melt pools of Nicolaides \citep{Nicolaides2026} and the brine pockets of Castlevetro
\citep{Castlevetro2026} form a complementary, largely equatorial class of liquid-water habitat that the
map flags through geomorphological diversity (as at Selk). The organic and methane-cycle features
driving the scores are supplied by the atmosphere Islam Rahim \citep{IslamRahim2026} shows to be
variable, while the deep energy budget is set by the interior Crick \citep{Crick2026} models. The
surface map thus points to \emph{where} the processes described by the rest of the group are most likely
to matter.

\begin{figure}[htbp]
  \centering
  \includegraphics[width=\linewidth]{habitability_map.png}
  \caption{Present-epoch (Cassini era) posterior mean surface habitability map, $P(H)$, from the
    Bayesian framework of \citet{Meadows2026}. North-polar lake shores (top) attain the highest
    habitability probability; the radar-bright, water-ice-rich terrain of Xanadu and the mountain
    ranges score lowest. Selk crater (equatorial), the Dragonfly landing site, scores only
    modestly on this present-day solvent metric; the map instead endorses the mission's choice via
    Selk's high geomorphological-diversity feature.}
  \label{fig:habmap}
\end{figure}
